\chapter{Architecture}\label{ch:architecture}

This part explains how the engine works, from the model's layers to the measurement method. This first chapter is the tour: the model, the byte budget that sets the speed limit, the decode loop, and a map of the code.

\section{The model}\index{Qwen3.8-27B}

\code{engine/config.py} reads the checkpoint's \code{config.json} and nothing else. The text model has 64 decoder layers, a hidden size of 5,120, a dense SiLU MLP of width 17,408 in every layer, and a vocabulary of 248,320 tokens. It has no mixture of experts. A 0.92~GB vision tower sits in the same checkpoint; the language-model loader skips it by name, and \code{engine/vision.py} loads it on its own for image input.

Every fourth layer (3, 7, 11 and so on up to 63) is full attention, 16 layers in all. Gated DeltaNet\index{Gated DeltaNet} fills the other 48. It is a recurrent layer: it keeps a fixed-size state per head instead of a KV cache that grows with every token.

\begin{figure}[htbp]
\centering
\begin{tikzpicture}[x=2.3mm, y=1mm]
  \foreach \i in {0,...,63} {
    \pgfmathtruncatemacro{\att}{mod(\i,4)==3 ? 1 : 0}
    \ifnum\att=1
      \fill[signal] (\i,0) rectangle +(0.82,9);
    \else
      \fill[cobalt!70] (\i,0) rectangle +(0.82,9);
    \fi
  }
  \foreach \i in {0,15,31,47,63} \node[lbl, below] at (\i+0.4,0) {\i};
  \node[lbl, anchor=west] at (0,13) {\textcolor{cobalt}{\rule{2.5mm}{2.5mm}} Gated DeltaNet, 48 layers: a 128$\times$128 fp32 state per head};
  \node[lbl, anchor=west] at (36,13) {\textcolor{signal}{\rule{2.5mm}{2.5mm}} full attention, 16 layers: KV rows};
\end{tikzpicture}
\caption[The layer pattern of Qwen3.8-27B]{The layer pattern of Qwen3.8-27B: three Gated DeltaNet layers, then one full-attention layer, sixteen times.}\label{fig:layers}
\end{figure}

\begin{table}[htbp]
\centering\small
\caption{The two layer kinds and the state each keeps per sequence (\arith, from the configuration's sizes).}\label{tab:layer-kinds}
\begin{tabularx}{\linewidth}{@{}L{.22\linewidth}rY@{}}
\toprule
\hd{layer kind} & \hd{count} & \hd{state per sequence} \\
\midrule
Gated DeltaNet & 48 & an fp32 state of 48 heads, each 128 by 128 (3.15 MB), and a 4-wide convolution tail (80 KB), per layer \\
full attention & 16 & BF16 keys and values, 4 KV heads of 256: 4 KB a token per layer \\
\bottomrule
\end{tabularx}
\end{table}

That split shapes most of the code. Published work on speculation and caching assumes attention, where undoing a token means moving a pointer back by one row. A recurrent state has absorbed every token and has no inverse. The engine needs exact ways to undo it, to branch it for a tree of guesses, and to restore it from a cache. Chapter~\ref{ch:speculation} shows the undo and the branch, and Chapter~\ref{ch:caches} the restore.

\section{The byte budget}\index{byte budget}

A decode step reads every weight once. Table~\ref{tab:bytes} sums the bytes from the checkpoint's safetensors headers.

\begin{table}[htbp]
\centering\small
\caption[The bytes one decode step reads, as shipped]{The bytes one decode step reads, as the vendor ships the model (\arith, from the safetensors headers).}\label{tab:bytes}
\begin{tabular}{@{}lr@{}}
\toprule
\hd{group} & \hd{bytes a step} \\
\midrule
48 Gated DeltaNet layers & 18.43 GB \\
16 attention layers & 5.96 GB \\
the output head, BF16, 248,320 rows by 5,120 & 2.54 GB \\
total, FP8 as shipped & 26.93 GB \\
\bottomrule
\end{tabular}
\end{table}

The three MLP projections are 63.6~\% of that. Quantising them first, then the rest, moved the budget down in steps, and each step had to pass the quality gate (Chapter~\ref{ch:quantisation}). Figure~\ref{fig:byte-budget} shows the steps and the floor each puts under a step at the board's paper bandwidth of 273~GB/s.

\begin{figure}[htbp]
\centering
\begin{tikzpicture}
\begin{axis}[tandem, ybar, height=6cm, width=.92\linewidth, bar width=22pt,
    ymin=0, ymax=31, ylabel={GB read per decode step},
    symbolic x coords={fp8,mlp,head,all}, xtick={fp8,mlp,head,all}, bar shift=0pt,
    xticklabels={{FP8 as shipped},{NVFP4 MLPs},{+ FP8 head},{+ NVFP4 GDN, attention}},
    x tick label style={align=center, text width=2.6cm},
    nodes near coords, nodes near coords style={font=\sffamily\scriptsize, align=center},
    point meta=explicit symbolic, enlarge x limits=0.18, grid=none]
  \addplot[fill=cobaltlight, draw=cobalt] coordinates
    {(fp8,26.93) [26.93 GB\\98.6 ms] (mlp,19.44) [19.44 GB\\71.2 ms] (head,18.17) [18.17 GB\\66.5 ms]};
  \addplot[fill=signal, draw=signal, bar shift=0pt] coordinates {(all,15.01) [15.01 GB\\55.0 ms]};
\end{axis}
\end{tikzpicture}
\caption[Bytes a step and the floor they set]{Bytes a step for each weight set, and the floor each sets under one step at 273~GB/s (\arith). The served set is the last bar. The board reads about 240~GB/s in practice, so the real floors are about 14~\% higher.}\label{fig:byte-budget}
\end{figure}

State is small beside the weights. The 48 recurrent states\index{recurrent state} are 151~MB together. The KV rows\index{KV rows} cost 64~KB a token across the 16 attention layers: 2.1~GB at 32k tokens and about 17~GB at the full 262,144-token context.

\section{The decode loop}\index{decode loop}\index{tokens per round}

Speed is one fraction:
\[
\text{tok/s} = \frac{\text{tokens committed a round}}{\text{draft ms} + \text{verify ms} + \text{commit ms}}.
\]
A round starts from the last committed token. The drafters propose a tree of guesses after it. One forward pass runs over the whole tree through all 64 layers. The engine reads the model's own choice at every node and keeps the longest path where the model agrees with the guesses, plus one more token, the model's own choice after the last accepted guess, so a round yields at least one token. Then it commits the state of that path and starts again.

On the served profile a round on the benchmark row takes 93.4~ms and commits 4.58 tokens, and the row's rate is 49.89 tok/s (\measured, the paper). The rate is the mean of the 50 requests' own rates, not the quotient of the two means (4.58 tokens in 93.4~ms is 49.0 tok/s). Without speculation the same box writes 13.69 tok/s at 72.8~ms a token. Chapter~\ref{ch:speculation} follows a round step by step.

\section{The parts}

Table~\ref{tab:code-map} maps each part of the engine to the files that implement it and the chapter that explains it.

\begin{table}[htbp]
\centering\small
\caption{A map of the code.}\label{tab:code-map}
\begin{tabularx}{\linewidth}{@{}L{.3\linewidth}YL{.14\linewidth}@{}}
\toprule
\hd{part} & \hd{where} & \hd{chapter} \\
\midrule
checkpoint layout and loading & \code{engine/loader.py}, \code{engine/config.py} & \ref{ch:adding-model} \\
the quantiser and its gate & \code{tools/quant_nvfp4.py}, \code{tools/quant_head.py}, \code{tools/quality_gate.py} & \ref{ch:quantisation} \\
matrix-product kernels & \code{tools/nvfp4_*.py}, \code{tools/fp8_linear.py}, \code{tools/head_gemv.py} & \ref{ch:kernels} \\
recurrent-layer kernels & \code{tools/gdn_*_kernels.py}, \code{engine/gdn.py} & \ref{ch:kernels} \\
target model, verify and commit & \code{engine/model.py}, \code{engine/tree.py} & \ref{ch:speculation} \\
drafters & \code{engine/drafters/} & \ref{ch:speculation} \\
routers: StairCut, length router, tree & \code{engine/lenrouter.py}, \code{engine/router.py}, \code{engine/prices.py} & \ref{ch:speculation} \\
sampling and penalties & \code{engine/sample.py}, \code{engine/penalty.py} & \ref{ch:exactness} \\
structured outputs & \code{engine/grammar.py} & \ref{ch:serving} \\
caches & \code{engine/cache.py}, \code{engine/disk_cache.py}, \code{engine/spill.py} & \ref{ch:caches} \\
settings registry & \code{engine/settings.py} & \ref{ch:operations} \\
the server & \code{server/app.py}, \code{server/stream.py}, \code{server/toolcall.py} & \ref{ch:serving} \\
the conductor & \code{server/conductor/} & \ref{ch:conductor} \\
live activity, metrics, ledger & \code{server/live.py}, \code{server/activity.py}, \code{server/metrics.py}, \code{server/ledger.py} & \ref{ch:dashboard} \\
the dashboard & \code{dashboard/}, \code{docs/contract/dashboard-v1/} & \ref{ch:dashboard} \\
the proxy and the collector & \code{proxy/} & \ref{ch:other-engines} \\
service scripts and the gate & \code{ops/} & \ref{ch:operations} \\
measurement tools & \code{tools/row3.py}, \code{tools/block_ab.py}, \code{tools/verify_spec.py} & \ref{ch:measurement} \\
\bottomrule
\end{tabularx}
\end{table}

\section{One sequence, or several}

The numbers in this part describe the loop of one sequence, which is what the shipped profile runs. With \code{--parallel N} the engine holds N sequences over the one copy of the weights, and requests that decode at the same time verify their trees in one forward pass, so the weights are read once for all of them (Section~\ref{sec:parallel-requests}). With one sequence, requests wait for the engine's lock in a bounded queue, the side lane serves a short request inside the running one (Section~\ref{sec:side-lane}), and the proxy relays other backends' requests without touching the engine's lock (Chapter~\ref{ch:install-dash}).
